% Standalone source-record documentation; no research attempts.
\documentclass[11pt]{article}
\usepackage[margin=25mm]{geometry}
\usepackage{fontspec,unicode-math}
\setmainfont{Latin Modern Roman}
\setmathfont{Latin Modern Math}
\usepackage{amsmath,amssymb,mathtools,xurl,hyperref}
\hypersetup{colorlinks=true,urlcolor=blue,linkcolor=blue}
\setcounter{secnumdepth}{0}
\setlength{\parindent}{0pt}
\setlength{\parskip}{5pt}
\setlength{\emergencystretch}{3em}
\newcommand{\sref}[2]{\hyperlink{source:#1:#2}{[S#2]}}
\newcommand{\cataloguescope}[1]{\paragraph{Source record.}#1}
\newcommand{\floor}[1]{\left\lfloor#1\right\rfloor}
\begin{document}
% UnsolvedMath ID 20002204; source catalogue code AIM-INFRASTRUCTURE-0104
\setcounter{section}{20002203}
\section{Breach-centered sampling of IBL classroom culture}\label{20002204}
{\small\sffamily UnsolvedMath ID 20002204 · Computer Science}
\subsection{Definitions and mathematical statement}

\paragraph{Shared notation.}Standard number systems are $\mathbb N,\mathbb Z,\mathbb Q,\mathbb R,\mathbb C$. Specialized notation and index ranges follow the cited source and the definitions below.


Inquiry-based learning (IBL) engages students in mathematical investigation and justification. A real-analysis course studies topics such as limits, continuity, and proofs from definitions. A classroom culture consists of expectations and practices for participation, explanation, feedback, and responsibility. A qualitative study analyzes observations and interviews to describe classroom practices and participants' interpretations. A case study examines a bounded classroom or course context; a comparison requires consistent categories and evidence collection across cases.
\paragraph{Record scope and qualifications.}
The catalogue title is retained as metadata. Named certificates, contracts, bounds or frameworks in that title are not asserted as results of the original source extract.
The extract includes proposed studies and appended discussion. Its empirical hypotheses are not established mathematical theorems, and no fixed causal identification model or validated coding scheme is supplied.

Catalogue screening limitation: Empirical education study and classroom-culture questions, rather than a mathematical statement.

The source snapshot labels this record partially solved; this is catalogue metadata rather than independent verification of a mathematical result or the current completion of a project.
The source text below is rendered literally where necessary; damaged formulas and composite extracts are not silently repaired into a new mathematical claim.
\paragraph{Statement recorded in the supplied source.}
(1) What strategies do IBL instructors use to create and sustain an appropriate classroom culture in IBL real analysis courses? (2) What are the essential aspects of appropriate classroom culture that faculty make explicit moves to establish and maintain, regardless of classroom structure or strategy for doing so?

Abstract: Because IBL instruction departs from the traditional (lecture-style) mathemat-ics classroom contract between teacher and students (and among peers), we anticipate that all IBL instructors use various strategies for creating and maintaining a classroom culture conducive to the expectations and goals of IBL pedagogy. While discussion and self-report of experienced IBL instructors suggests that these strategies are diverse, we hypothesize that the diverse strategies are used to accomplish a more universal set of goals. We pro-pose a study of these teaching moves, their purposes, and their efficacy through qualitative observation of IBL classrooms. We anticipate doing so by identifying key points in which faculty make explicit moves to create and sustain a desirable classroom culture before con-ducting faculty and student interviews around video clips of such events. Eliciting student interpretation of such events and how they arise to breaches of the intended classroom ex-pectations, we anticipate being able to learn about students understanding and adoption of these expectations for learning.

• Video record early classes and then a sampling later in the semester.

• Attend to initial solicitation for buy in to IBL and then to breaches of classroom contract that are negotiated. 3

• Conduct member check interviews for instructor and student interpretations of class-room expectations and the justification for/sources of such expectations. It was proposed to measure student perceptions twice per semester, possibly after 4 and 12 weeks (based upon IBL instructors experiences of the recurrent progression of students affective responses in such courses).

• Possible student questioning strategy: If you could choose how classroom time is spent in the mathematics course you are taking, what percentage of the time would you spend on: professor lecture, student presentations, group work on challenging problems, etc?

• We assume many students will come into an IBL real analysis course with a baseline set of expectations informed by traditional mathematics instruction in which students attempt tasks after receiving relevant direct instruction, students are assessed on private work without publishing student work to other students, and the professor and textbook are the primary sources of new mathematics in the class. IBL is taken to entail the expectation that students attempt proof tasks for which they have not received direct instruction, students present their own work to the class for evaluation, and that students author new mathematical content into the classroom.

• Relevant categories we anticipate may arise in the study:

- Legitimizing failure (or linking success to hard work and repeated attempts)

- Soliciting buy in to the alternative classroom expectations

- Negotiating classroom expectations (of students and teacher)

- Endorsing standards for acceptable proof

- Curtailing undesirable mathematical practices

• We anticipate observing IBL instructors of various levels of experience. Novice in-structors might provide more opportunity to identify breaches and difficulties that require intervention. Observing experienced IBL instructors will aide in identifying successful strategies, since documenting and publishing such strategies is an impor-tant research product.

Learning through proof presentations. The overarching goal of this investigation is to under-stand what students learn in IBL courses when listening to peer proof presentations. Two relevant hypotheses were identified:

Hypothesis 1: Students learn more from proof presentations in an IBL course because they have already attempted the proof tasks before seeing the completed solution.

Hypothesis 2: Students learn more from proof presentations in an IBL course because they pay closer attention to find and identify mistakes, which may depend upon how student listening is guided by the instructor. We propose two possibly independent studies for these two hypotheses.

Research Questions: \sref{20002204}{2}
\paragraph{Source status.}
The supplied catalogue labels this source record partially solved; that metadata was checked against the supplied snapshot on 2026-10-11. The label is preserved for provenance and is not an independent verification of present open status or mathematical correctness.
\subsection{Short English statement}
Study how IBL instructors establish classroom culture and how students experience it.
\subsection{Sources}
\begin{itemize}
\item[\hypertarget{source:20002204:1}{S1}] Ulam AI and the UnsolvedMath Contributors, UnsolvedMath: A Curated Collection of Open Mathematics Problems, record 20002204 (AIM-INFRASTRUCTURE-0104). Catalogue curation and metadata: CC BY 4.0. The source extract is retained; definitions, English summaries and qualifications are editorial additions.
\url{https://huggingface.co/datasets/ulamai/UnsolvedMath}.
\item[\hypertarget{source:20002204:2}{S2}] AIM Workshop Problem Lists, source linked by catalogue record AIM-INFRASTRUCTURE-0104.
\url{http://aimath.org/pastworkshops/iblanalysisrep.pdf}.
\end{itemize}
\paragraph{Preparation and review.}
The original source and UnsolvedMath contributors retain attribution for the source extract. Codex (GPT-6; exact serving revision not exposed) assisted with definitions, English summaries, scope notes and formatting. No new human mathematical authorship or independently verified solution is asserted. Permanent source ID 20002204 and code AIM-INFRASTRUCTURE-0104 checked against the complete supplied registry and current repository. Distinct upstream records, including repeated questions and broad or damaged extracts, are retained for complete coverage of the source CS category. A repeated record is not a new independent conjecture. Complete mathematical-equivalence and current-literature review has not been performed. Preparation and final review follow contribution prompts C and D, with source gaps disclosed. Structural and rendering checks do not establish a new theorem or open status.
\label{end:20002204}
\end{document}
